% ===================================================================
%  اللوحة رقم ٧ — الضيعة بالشتوية. --- بيت المزرعة بالربيع.
%
%  Traduction, faite en 2026, en arabe levantin d'aujourd'hui, sur
%  l'ido de texto/io. Regles de la colonne : voir 10-lawha-01.tex.
%  Ouverture de la DEUXIEME SERIE : « السلسلة التانية ». Deux scenes,
%  deux numerotations. Un bloc marque « ¶2 » par ordo.py,
%  t07-c1-07-1 : une seule cle, deux alineas.
%
%  LA CHATAIGNE (42) EST « كستنا », ET IL FAUT LIRE CE RENVOI A COTE
%  DE CELUI DE LA COLONNE GUJARATIE. Au meme numero, dans le meme
%  alinea, le gujarati a REFUSE : il avait « શેકેલી સિંગ »,
%  l'arachide grillee qu'on vend en cornet de papier a Ahmedabad
%  l'hiver — meme scene, meme papier, meme saison — et il a du ecrire
%  « ચેસ્ટનટ » parce que la chose n'etait pas la meme. Ici il n'y a
%  rien a refuser : le chataignier pousse sur le Mont-Liban et dans
%  le Qalamoun, on vend les « كستنا » grillees au coin des rues de
%  Beyrouth et de Damas en janvier, dans le meme cornet, et l'enfant
%  du livret se rechauffe les mains avec exactement ce que se
%  rechaufferait un enfant d'ici. DEUX COLONNES, UN RENVOI, DEUX
%  REPONSES OPPOSEES — et c'est la geographie qui decide, pas la
%  langue.
%
%  TOUT CE TABLEAU D'HIVER EST DANS CE CAS. La luge (32), les patins
%  (26), le bonhomme de neige (39), les boules de neige (38), la
%  glace de la riviere (25) : rien a transposer et rien a decrire. Le
%  Levant a des montagnes qui blanchissent quatre mois par an, et la
%  colonne gujaratie avait passe son tableau 7 a expliquer un hiver
%  qu'elle ne pouvait pas montrer.
%
%  LES METIERS DU VILLAGE CONTINUENT LE TABLEAU 5. Le cordonnier (7)
%  est « الكندرجي », de kundura, la chaussure turque, plus le suffixe
%  -ci ; l'aubergiste (4) est « الخانجي », de خان, le caravanserail.
%  Le charron (9) et le voiturier gardent « عربجي ». TROIS TABLEAUX
%  DE SUITE, LE MEME SUFFIXE, ET IL N'A PAS FINI.
%
%  Et le porc (77, 78) s'ecrit sans note, comme partout ailleurs dans
%  le releve — mais pour une raison qui n'avait pas encore servi. Au
%  Gujarat et chez les Haoussa, l'interdit valait pour presque tout le
%  monde et le mot survivait quand meme. Ici, LE VILLAGE EST MIXTE :
%  les chretiens du Liban et de Syrie en elevent et en mangent, leurs
%  voisins non, et le mot « خنّوص » pour le porcelet est du levantin
%  parfaitement ordinaire parce que quelqu'un, dans la meme rue, s'en
%  sert. La loi tient toujours ; c'est sa raison qui change.
% ===================================================================

%%K t07-noto-1 noto
\VUnotes{143.2mm}{%
(*) شوف تفاصيل الأكل باللوحتين ٦ و١٣.%
}

%%K t07-apar-1 apar
\VUsaut{4.0mm}
\VUcentre{13.2pt}{90}{السلسلة التانية}
\VUsaut{3.0mm}
\VUfilet{16mm}
\VUsaut{5.6mm}
\VUtitre{15.0pt}{60}{اللوحة رقم ٧}
\VUsaut{3.0mm}

%%K t07-tit-1 sub
\VUcentre{12.0pt}{0}{\textit{المشهد الأول.}}
\VUsaut{3.0mm}

%%K t07-tit-2 sub
\VUpk{10.2pt}{40}{الضيعة بالشتوية.}
\VUsaut{4.0mm}

%%K t07-c1-01-1 p
1. --- بعطلة رأس السنة، رحت مع أبي على نوفوربو، ضيعة زغيرة كان لازم
يخلّص فيها كم شغلة تجارية.

%%K t07-c1-02-1 p
2. --- أخدنا \VUgras{عربة البوسطة} \textsuperscript{(11)} اللي بتمشي بين
المدينة وهالضيعة ؛ بس لأنه كان في برد كتير، هالسفرة بعربايّة مش كتير
مريحة ما كانت حلوة.

%%K t07-c1-03-1 p
3. --- لهيك انبسطنا كتير لمّا شفنا \VUgras{العربجي} عم يوقّف عربايته
عالساحة، قدّام \VUgras{خان} \textsuperscript{(2)} \textit{الدبّ الأبيض}.
و\VUgras{الخانجي} \textsuperscript{(4)} كان مستنّينا عند عتبة بيته، عم يدخّن
\VUgras{غليونه} بهدوء.

%%K t07-c1-03-2 p
عزمنا نفوت وأنا ما استنّيت حدا يترجّاني لقعد قدّام نار الحطب اللي كانت
عم تفرقع بالموقدة. وهونيك تمسخرت كتير على \VUgras{هوا الشمال} القارص
اللي عم يصرصر \VUgras{اللافتة} \textsuperscript{(3)} العتيقة وعم يشيل
بدوّامات سودا \VUgras{الدخان} \textsuperscript{(1)} الطالع من المدخنة.

%%K t07-c1-04-1 p
4. --- كان وقت الغدا. وبلا ما نستنّى أكتر، أكلنا منيح الأكل البسيط بس
الطيّب اللي قدّموه إلنا (*).

%%K t07-tit-3 sub
\VUpk{10.2pt}{40}{كم زيارة.}
\VUsaut{3.0mm}

%%K t07-c1-05-1 p
5. --- بعد الغدا، رحت مع أبي، اللي هوّي وكيل تأمين، على زباينه. أول شي
زرنا \VUgras{الحدّاد} \textsuperscript{(5)}، اللي ساكن جنب الخان. كان مشغول
عم ينعّل حصان. وأنا انبهرت بقوّة مساعده اللي مسّك منيح على مريوله
الجلد \VUgras{حافر} الحصان، وقت الحدّاد عم يثبّت \VUgras{النعل}
\textsuperscript{(6)} بضربات مطرقة قوية.

%%K t07-c1-06-1 p
6. --- وبعدين رحنا على \VUgras{الكندرجي} \textsuperscript{(7)} تبع الضيعة،
يعقوب العتيق. مع إنّه إلو لافتة \VUgras{جزمة} كتير حلوة، اكتفى
إنّه يحطّ نعال جديدة لفردة كنادر عتيقة مخروقة.

%%K t07-c1-06-2 p
وقلّنا الختيار وهوّي عم يضحك : « بالمدينة، كنت « صانع كنادر » وما كان
عندي زباين. وهون، أنا « مصلّح كنادر » والشغل كتير. »

%%K t07-c1-07-1 p
7. --- وحكانا عن جاره \VUgras{الحلّاق} \textsuperscript{(8)}، اللي زباينه
مش راضيين. \VUgras{أمواسه} دايماً مثلومة و\VUgras{مقصّاته} أبداً ما
بتقصّ منيح.

بس لأنه هوّي الحلّاق الوحيد بالضيعة، مضطرّين طبعاً يحلقوا عنده ويقصّوا
شعرهن عنده. وبعدين هوّي رجّال بخيل ومش لطيف.

%%K t07-c1-08-1 p
8. --- ولحسن الحظّ، باقي \VUgras{الجيران} مش متلو. مثلاً \VUgras{صانع
العربايات} \textsuperscript{(9)} رجّال طيّب، شغّيل وبيحبّ يساعد. متعة
تصلّح عربايتك عنده. شغله دايماً بيكون خالص.

%%K t07-c1-09-1 p
9. --- ويعقوب العتيق كمان ما اشتكى من \VUgras{السرّاج} \textsuperscript{(10)}،
اللي بعض المرّات بيساعده يخيّط \VUgras{عدّة الخيل} لمّا يكون الشغل
مستعجل.

%%K t07-c1-09-2 p
وأبي سأله عن عيلته : « الكل منيحين. حفيدتي بـ\VUgras{المدرسة}
\textsuperscript{(12)}. و\VUgras{معلّمتها} \textsuperscript{(13)} كتير راضية
عنها، هيّي \VUgras{تلميذة} \textsuperscript{(14)} منيحة. مش متل ابني الشقي ؛
هوّي ما بيفكّر غير باللعب. و\VUgras{المعلّم} \textsuperscript{(15)} ما عم
يقدر يعلّمه شي. »

%%K t07-tit-4 sub
\VUsaut{3.0mm}
\VUpk{10.2pt}{40}{حكي الضيعة.}
\VUsaut{3.0mm}

%%K t07-c1-10-1 p
10. --- وبعدين الختيار قصّ إلنا أخبار الضيعة. رح يبنوا مدرسة جديدة،
لأنه اللي بـ\VUgras{دار البلدية} صارت زغيرة كتير. وهاد سهل، لأنه بيكفي
يشتروا قسم من جنينة \VUgras{الطحّان}. وهالشي رح يفيد هالرجّال المسكين
كتير. وبسبب البرد، \VUgras{حجار الطحن} تبع \VUgras{الطاحون}
\textsuperscript{(16)} ما عادت تقدر تطحن القمح، لأنه \VUgras{الدولاب}
\textsuperscript{(17)} وقّف مجلّد من \VUgras{جليد} \textsuperscript{(25)}
\VUgras{الساقية} \textsuperscript{(23)}.

%%K t07-c1-11-1 p
11. --- « وبخصوص البنا، شفتوا \VUgras{كنيستنا} \textsuperscript{(18)}
الجديدة ؟ كتير حلوة تحت غطاها التلج. و\VUgras{الغربان}
\textsuperscript{(19)}، اللي ما عادوا يعرفوا برج أجراسهن العتيق، عم
يحطّوا بحزن على شجرة \VUgras{المقبرة} \textsuperscript{(20)} الكبيرة. »

%%K t07-c1-12-1 p
12. --- وفكرة المقبرة ذكّرت الكندرجي بالناس اللي كان أبي بيعرفهن واللي
ماتوا السنة الماضية. « قدّيش \VUgras{قبر} \textsuperscript{(21)}، يا
سيّدي، انضافوا عالباقيين، تحت \VUgras{السرو} \textsuperscript{(22)} ! الموت
حصد كاتبنا بالعدل العتيق. وكل الضيعة حضروا \VUgras{جنازته}. »

%%K t07-c1-13-1 p
13. --- « هيّو خليفته عم يتزحلق عالساقية. هلّق كان جايي عندي ليصلّح
رباط \VUgras{الباتيناج} \textsuperscript{(26)} تبعه اللي انقطع. »

%%K t07-c1-14-1 p
14. --- « وهيّو كمان على ضفّة الساقية، \VUgras{خفير الضيعة}
\textsuperscript{(27)} الجديد. هوّي درك سابق وبيخوّف أشقيا الضيعة. عم
يهربوا أول ما يشوفوا \VUgras{طماقاته} \textsuperscript{(29)}
و\VUgras{قبّعته العسكرية} \textsuperscript{(28)}. »

%%K t07-c1-15-1 p
15. --- « آه ! هيّو يوسف العتيق عم يسوق \VUgras{عربايّة حزم الحطب}
\textsuperscript{(30)} لصاحب الفرن. --- صحيح، قال أبي، عم عرف
\VUgras{بغاله} \textsuperscript{(31)} المجنّزرة. أنا لازم حكي معه، رح
استفيد من الفرصة. لبكرا يا أبو يعقوب. »

%%K t07-c1-16-1 p
16. --- وأول ما طلعنا، ولاد الضيعة طلعوا من المدرسة وركضوا عالـ\VUgras{مزلقان}
\textsuperscript{(33)}، ومخاطرين يوقعوا ويكسروا إجر أو إيد
بـ\VUgras{الوقعة} \textsuperscript{(34)}. وواحد منهن أخد \VUgras{زحّافته}
\textsuperscript{(32)} من صانع العربايات، وقعّد فيها واحد من رفقاته،
وراح يركض.

%%K t07-c1-17-1 p
17. --- وغيرهن ركضوا على \VUgras{كومة التلج} \textsuperscript{(37)} اللي
جنب \VUgras{الجسر} \textsuperscript{(40)} وبلّشوا يقصفوا \VUgras{رجل التلج}
\textsuperscript{(39)} اللي عملوه مبارح واللي حطّوا عليه برنيطة وغليون
وشنتا مدرسة عالكتف.

%%K t07-c1-18-1 p
18. --- وما بيهمّهن كتير الهوا المجلّد ولا البرد. وأكترهن ولا معهن
\VUgras{لفحة} \textsuperscript{(35)}، وليدفوا حالهن بعد معركة \VUgras{كور
التلج} \textsuperscript{(38)}، بيكفيهن كفّ \VUgras{كستنا}
\textsuperscript{(42)} كتير سخنة بيشتروها من \VUgras{البيّاعة} العتيقة
ياكوبينا، اللي عم ترجف من البرد تحت \VUgras{شالها}
\textsuperscript{(43)} الرفيع كتير.

%%K t07-tit-5 sub
\VUsaut{3.0mm}
\VUcentre{12.0pt}{0}{\textit{المشهد التاني.}}
\VUsaut{3.0mm}

%%K t07-tit-6 sub
\VUpk{10.2pt}{40}{بيت المزرعة بالربيع.}
\VUsaut{4.0mm}

%%K t07-c2-01-1 p
1. --- مع كل ملذّات الشتوية، منسلّم بفرح عميق على رجعة
\VUgras{الربيع}.

%%K t07-c2-01-2 p
شو ما أحلى المنظر اللي بيبيّنه حوش مزرعة، عالمسا، بشهر أيار !

%%K t07-c2-02-1 p
2. --- وعند الأفق \VUgras{الشمس} \textsuperscript{(1)} عم تغيب على مهلها،
وعم تغرق \VUgras{الضيعة} \textsuperscript{(3)} بـ\VUgras{أشعّتها}
\textsuperscript{(2)} الأرجوانية. و\VUgras{الفلّاح} \textsuperscript{(6)}
النشيط عم يستعجل يفلح \VUgras{حقله} \textsuperscript{(4)}.

%%K t07-c2-02-2 p
وبـ\VUgras{محراثه} \textsuperscript{(5)} المجنّزر بـ\VUgras{حصان}
\textsuperscript{(7)} قوي، عم يشقّ \VUgras{التلم} \textsuperscript{(8)} اللي
فيه \VUgras{الزرّاع} \textsuperscript{(9)} رح يرمي بكفّه المليانة
\VUgras{البذار} اللي رح يعطي السنابل الدهبية.

%%K t07-c2-03-1 p
3. --- و\VUgras{الحصّادين} \textsuperscript{(11)} بمناجلهن حصدوا عشب
المرج، واللي بيقلّبوا القش نشّفوا \VUgras{القش} \textsuperscript{(10)} وهنّي
عم يقلّبوه بـ\VUgras{المذاري} \textsuperscript{(12)} والمجارف، وهيّنن
راجعين.

%%K t07-c2-03-2 p
وبعضهن عم يمشوا قدّام العربايّة التقيلة اللي عم يجرّها البقر ؛ حاملين
عدّتهن عالكتف. وغيرهن، أكسل، قعدوا مرتاحين على القش المعطّر.

%%K t07-c2-04-1 p
4. --- ووهنّي مارقين عم يسلّموا بودّ على \VUgras{الجنايني}
\textsuperscript{(16)} المشغول بـ\VUgras{البستان} \textsuperscript{(13)} عم
يقلّم \VUgras{شجر الفواكه} ويطعّم \VUgras{شجر الإجّاص}
\textsuperscript{(14)}. و\VUgras{شجر الخوخ} \textsuperscript{(15)} عم يزهّر،
وبـ\VUgras{جنينة الخضرة} \textsuperscript{(17)} المزبّلة منيح، الخضرة
و\VUgras{الملفوف} و\VUgras{الخسّ} صاروا منيحين من هلّق.

%%K t07-c2-04-2 p
وعالحيط الواطي، \VUgras{طاووس} \textsuperscript{(18)} حلو عم يفتح
\VUgras{ديله}، ليفرجي \VUgras{ريشه} الحلو كتير.

%%K t07-tit-7 sub
\VUpk{10.2pt}{40}{حوش المزرعة.}
\VUsaut{3.0mm}

%%K t07-c2-05-1 p
5. --- أبواب \VUgras{الإسطبل} \textsuperscript{(19)} مفتوحة ومنشوف المعلف
و\VUgras{فرشة} \textsuperscript{(23)} \VUgras{الفرس} \textsuperscript{(21)} اللي
\VUgras{السايس} \textsuperscript{(20)} هلّق فكّها من الجنزير.
و\VUgras{المهر} \textsuperscript{(22)} المضحك كتير بإجريه الطويلة عم يلحق
إمّه وهوّي عم ينطّ.

%%K t07-c2-06-1 p
6. --- وجنب \VUgras{البير} \textsuperscript{(29)}، \VUgras{البقرات}
\textsuperscript{(25)} رايحات يشربوا من \VUgras{الحوض} \textsuperscript{(26)}
الخشب اللي بيخدمهن كمشرب. و\VUgras{العجل} \textsuperscript{(24)} عم يستفيد
من الفرصة ليرضع، و\VUgras{التور} \textsuperscript{(27)} وهوّي عم يشمّ
ريحة العلف المنيحة، عم يخور بفرح وعم يرفع بفخر راسه
بـ\VUgras{قرونه} \textsuperscript{(28)} القوية.

%%K t07-c2-07-1 p
7. --- والكل عم يشتغلوا. \VUgras{الخدّامة} \textsuperscript{(32)} ماسكة
بإيدها \VUgras{حبل} \textsuperscript{(31)} البير. عم تسحب ميّ
بـ\VUgras{السطل} \textsuperscript{(30)} لتسقي المواشي. وجنب \VUgras{البيدر}
\textsuperscript{(33)}، رجّال عم يجمع القش، وقدّام باب \VUgras{البيت}
\textsuperscript{(34)} \VUgras{غزّالة} \textsuperscript{(35)} عتيقة، بدولاب
غزلها وبـ\VUgras{مغزلها}، عم تغزل \VUgras{الكتّان} أو \VUgras{القنّب}
متل زمان.

%%K t07-tit-8 sub
\VUsaut{3.0mm}
\VUpk{10.2pt}{40}{صاحب المزرعة.}
\VUsaut{3.0mm}

%%K t07-c2-08-1 p
8. --- \VUgras{صاحب المزرعة} \textsuperscript{(36)} عم يدير كل
هالأشغال وعم يعطي تعليمات لعمّاله. عم يوصّيهن يحطّوا تحت السقف
\VUgras{المشط} \textsuperscript{(40)} و\VUgras{الكزمة} \textsuperscript{(39)}
اللي نسيوهن جنب \VUgras{بيت} \textsuperscript{(38)} \VUgras{الكلب}
\textsuperscript{(37)}.

%%K t07-c2-09-1 p
9. --- وصياحه بيخوّف \VUgras{حمامة} \textsuperscript{(41)} عم تطير بكل
جناحيها لـ\VUgras{برج الحمام} \textsuperscript{(42)}، و\VUgras{الدجاجات}
\textsuperscript{(43)} اللي عم يستعجلوا يرجعوا عالـ\VUgras{خمّ}
\textsuperscript{(44)}.

%%K t07-c2-10-1 p
10. --- وقدّام \VUgras{الزريبة} \textsuperscript{(48)}، هيّو
\VUgras{الراعي} \textsuperscript{(45)} العتيق عم يرجّع \VUgras{قطيعه}
\textsuperscript{(49)}. وماسك \VUgras{عصاه} \textsuperscript{(46)}، اللي
أبداً ما بتفارقه، متل \VUgras{عبايته} \textsuperscript{(47)} المبهتة، عم
يسوق عالحظيرة \VUgras{خرفان} \textsuperscript{(50)}
و\VUgras{نعاج} \textsuperscript{(53)} و\VUgras{حملان} \textsuperscript{(54)}
و\VUgras{معزات} \textsuperscript{(51)} و\VUgras{جديان}
\textsuperscript{(52)}. و\VUgras{كلبه} \textsuperscript{(55)} الوفي أرغوس
بيجي بالآخر وبيحثّ بنبحه المتأخّرين.

%%K t07-c2-11-1 p
11. --- وكل هالضجّة وهالحركة ما عم تزعزع هدوء \VUgras{البقرة الحلوب}
\textsuperscript{(56)} الطيّبة اللي عم يرنّ \VUgras{جرسها}
\textsuperscript{(57)}، وقت عم يحلبوها قدّام باب \VUgras{الحظيرة}
\textsuperscript{(58)}.

%%K t07-c2-12-1 p
12. مبيّن إنّها فاهمة إنّه لازم تعطي حليبها، لتقدر \VUgras{صاحبة
المزرعة} \textsuperscript{(60)} تحضّر \VUgras{الزبدة} بـ\VUgras{المخيض}
\textsuperscript{(61)} أو \VUgras{الأجبان} \textsuperscript{(64)} الممتازة،
اللي عم تنقّط من اللوح المحطوط عالـ\VUgras{لقّان}
\textsuperscript{(63)}.

%%K t07-c2-13-1 p
13. --- وكل \VUgras{بيت الحليب} \textsuperscript{(59)} بيحافظ عليه كتير
نضيف \VUgras{عامل المزرعة} \textsuperscript{(65)} اللي هلّق غسل
\VUgras{أوعية الحليب} \textsuperscript{(62)} بالسطل الكبير اللي حامله
بإيده.

%%K t07-tit-9 sub
\VUsaut{3.0mm}
\VUpk{10.2pt}{40}{الخمّ.}
\VUsaut{3.0mm}

%%K t07-c2-14-1 p
14. --- وبقباقيبه الخشب التخينة، عم يمشي شوي بثقل، وهيك عم يهرّب
\VUgras{الديك الحبش} \textsuperscript{(68)}، و\VUgras{البطّات}
\textsuperscript{(66)}، و\VUgras{البطّ} \textsuperscript{(67)}
و\VUgras{الوزّ} \textsuperscript{(69)} اللي عم يفتحوا \VUgras{مناقيرهن}
\textsuperscript{(70)} عالآخر وعم يصيحوا بقلق.

%%K t07-c2-14-2 p
و\VUgras{الديك} \textsuperscript{(71)}، اللي أكتر شراسة، عم ينصب
\VUgras{عرفه} \textsuperscript{(72)} الأحمر، وعم ينفض ريش \VUgras{ديله}
\textsuperscript{(73)}، وعم يصيح « كوكوريكو » بصوت عالي.

%%K t07-c2-15-1 p
15. --- و\VUgras{الدجاجة الإم} \textsuperscript{(74)} عم تنقّر على
\VUgras{الزبل} \textsuperscript{(81)} مع \VUgras{صيصانها}
\textsuperscript{(75)}. و\VUgras{الخنّوص} \textsuperscript{(78)} عم يهرب
لعند إمّه، \VUgras{الخنزيرة} \textsuperscript{(77)} الضخمة اللي منشوفها
جنب زريبة الخنازير.

%%K t07-c2-16-1 p
16. --- وبأقفاصهن، \VUgras{الأرانب} \textsuperscript{(79)} عم ياكلوا
بشراهة \VUgras{العشب} \textsuperscript{(80)} الطريّ اللي بتجيبلهن بنت صاحب
المزرعة. ويوانينا كتير بتحبّ أرانبها، وكل يوم، بعد ما تاخد
\VUgras{البيض} \textsuperscript{(76)} الطازة من الخمّ، بتيجي تزور رفقاتها
الزغار.

\VUsaut{6.0mm}
\VUfilet{22mm}
